\documentclass[11pt,a4paper]{article}
\usepackage[utf8]{inputenc}
\usepackage[T1]{fontenc}
\usepackage{amsmath,amssymb,amsthm}
\usepackage{booktabs}
\usepackage{geometry}
\usepackage{hyperref}
\usepackage[table]{xcolor}
\usepackage{enumitem}
\usepackage{graphicx}

\geometry{margin=2.5cm}
\hypersetup{colorlinks=true,linkcolor=blue,citecolor=blue,urlcolor=blue}

\definecolor{cgreen}{HTML}{2ea043}
\definecolor{corange}{HTML}{d29922}
\definecolor{cred}{HTML}{cf222e}

\title{Identifiability of the Connectome Inverse Problem:\\
Equivalence Classes and Ease-of-Identifiability}
\date{}

\begin{document}
\maketitle

\section{The Neural Model}

The flyvis graded-voltage ODE for each postsynaptic neuron $i$ is
\begin{equation}
\tau_i\frac{dv_i(t)}{dt} = -v_i(t) + V_i^{\text{rest}}
+ \sum_{j\in\mathcal{N}_i} \mathbf{W}_{ij}\,
  \text{ReLU}\!\big(v_j(t)\big)
+ I_i(t)
\label{Eq:ODE}
\end{equation}
where $\tau_i$ and $V_i^{\mathrm{rest}}$ are cell-type parameters, $\mathbf{W}_{ij}$ is the
connectome-constrained synaptic weight from neuron $j$ to $i$, and $I_i(t)$ is the
external stimulus. \textbf{The inverse problem}: given observed trajectories
$\{v_i(t)\}_{t=0}^{T}$, stimuli $\{I_i(t)\}$, and the binary adjacency matrix
$\mathbf{A}$, recover the connectivity weights $\mathbf{W}$, time constants $\tau$,
and resting potentials $V^{\text{rest}}$.

Unlike the teacher-student setting of Martinelli et al.\ (2024), where the teacher
is an MLP with discrete permutation and sign-flip symmetries, here the ``teacher''
is a physical ODE simulator and the symmetry is \emph{continuous}: the forward map
$\mathcal{F}: \mathbf{W} \mapsto \{v_i(t)\}$ is many-to-one, with fibers that are
affine subspaces rather than finite point clouds.

%==============================================================================
\section{Equivalence Class of Zero-Rollout-Loss Connectomes}
%==============================================================================

\subsection{Rearranging to a linear system per neuron}

Isolating the incoming weights in Eq.~\ref{Eq:ODE}:
\begin{equation}
\sum_{j\in\mathcal{N}_i} \mathbf{W}_{ij}\, \text{ReLU}\!\big(v_j(t)\big)
= \tau_i \frac{dv_i(t)}{dt} + v_i(t) - V_i^{\text{rest}} - I_i(t)
\label{Eq:ODE2}
\end{equation}
Define the presynaptic activity $h_j(t) = \text{ReLU}(v_j(t))$ and the residual
$b_i(t) = \tau_i \dot{v}_i(t) + v_i(t) - V_i^{\text{rest}} - I_i(t)$.
Stacking $T$ time points yields the linear system:
\begin{equation}
\mathbf{H}_i\, \mathbf{w}_i = \mathbf{b}_i, \qquad
\mathbf{H}_i \in \mathbb{R}^{T \times d_i},\;
\mathbf{w}_i \in \mathbb{R}^{d_i},\;
\mathbf{b}_i \in \mathbb{R}^T
\end{equation}
The \textbf{equivalence class} of zero-rollout-loss connectomes for neuron $i$ is
the affine manifold:
\begin{equation}
[\mathbf{w}_i^*] = \mathbf{w}_i^* + \ker(\mathbf{H}_i)
\end{equation}
Any $\hat{\mathbf{w}}_i \in [\mathbf{w}_i^*]$ produces identical dynamics to the
true $\mathbf{w}_i^*$. This is a \emph{continuous} equivalence class of dimension
$\dim\ker(\mathbf{H}_i) = d_i - \text{rank}(\mathbf{H}_i)$, in contrast to the
discrete (finite) equivalence classes arising from permutation symmetry in
teacher-student MLP recovery.

%------------------------------------------------------------------------------
\subsection{Global SVD bound on the equivalence class dimension}
\label{subsec:global-svd}
%------------------------------------------------------------------------------

By rank-nullity:
\begin{equation}
\dim[\mathbf{w}_i^*] = d_i - \text{rank}(\mathbf{H}_i)
\end{equation}
We estimate $\text{rank}(\mathbf{H}_i)$ via SVD of the subsampled population
activity $\mathbf{H} \in \mathbb{R}^{8000 \times 982}$. The effective rank at
99\% cumulative variance is 45, giving:
\begin{equation}
\dim[\mathbf{w}_i^*] \approx \max(0,\; d_i - 45)
\end{equation}
Summing over all neurons:
\begin{equation}\label{eq:null}
\sum_{i=1}^{13,697} \max(0,\; d_i - 45) = 115{,}223
\quad (\sim 26.5\% \text{ of all edges lie in non-identifiable directions})
\end{equation}

\begin{center}
\begin{tabular}{rrrr}
\toprule
Variance threshold & Effective rank $r$ & $\dim[\mathbf{w}^*]$ & \% non-identifiable \\
\midrule
90\%  & 1   & 420,415 & 96.8\% \\
95\%  & 2   & 407,427 & 93.9\% \\
\textbf{99\%}  & \textbf{45}  & \textbf{115,223} & \textbf{26.5\%} \\
99.5\% & 90  & 26,413  & 6.1\%  \\
99.9\% & 238 & 0       & 0.0\%  \\
\bottomrule
\end{tabular}
\end{center}

%------------------------------------------------------------------------------
\subsection{Per-neuron identifiability}
%------------------------------------------------------------------------------

Computing the SVD of $\mathbf{H}_i$ separately for each postsynaptic neuron $i$,
a neuron is \textbf{identifiable} if $d_i \leq r_i$ (all incoming weights
constrained by dynamics) and \textbf{non-identifiable} if $d_i > r_i$ (some
weights free within $[\mathbf{w}_i^*]$). Per-neuron analysis reveals 80.3\% of
edges are non-identifiable at 99\% variance, stronger than the global bound
(26.5\%), reflecting neuron-level input correlations from duplicate-type groups
(\S\ref{subsec:duplicate}).

\begin{center}
\begin{tabular}{rrrr}
\toprule
Variance threshold & Mean rank $\bar{r}_i$ & $\dim[\mathbf{w}^*]$ & \% non-identifiable \\
\midrule
90\%  & 1.2 & 417,563 & 96.2\% \\
95\%  & 2.0 & 406,901 & 93.7\% \\
\textbf{99\%}  & \textbf{6.2}  & \textbf{348,545} & \textbf{80.3\%} \\
99.5\% & 8.9 & 312,396 & 72.0\% \\
99.9\% & 15.4 & 223,775 & 51.5\% \\
\bottomrule
\end{tabular}
\end{center}

%------------------------------------------------------------------------------
\subsection{Duplicate-type symmetry: the structural source of non-identifiability}
\label{subsec:duplicate}
%------------------------------------------------------------------------------

The columnar organization of the flyvis connectome (217 columns, each covering
one visual field position) forces same-type neurons across columns to produce
time-shifted copies of the same signal. In E\&C terminology, these constitute
\textbf{duplicate-type neuron groups}: when postsynaptic neuron $i$ receives
input from $k$ neurons of the same cell type, their $k$ activity columns in
$\mathbf{H}_i$ are strongly correlated, contributing rank $\approx 1$ instead of
$k$.

When activities are identical ($\mathbf{h}_{j_1} = \cdots = \mathbf{h}_{j_k}$), a
perturbation $\boldsymbol{\delta} \in \ker(\mathbf{H}_i)$ requires
$\sum_{m=1}^{k}\delta_{j_m} = 0$ --- a sum-zero constraint leaving $k-1$ free
directions, exactly as duplicate neurons in an overparameterised student can
redistribute output weights summing to the teacher activation. Summing over all
neurons and cell types:
\begin{equation}
\dim[\mathbf{W}^*] = \sum_{i=1}^{N} \sum_{\alpha:\,k_{i\alpha}>1} (k_{i\alpha}-1)
\approx 121{,}100 \quad (\sim 28\% \text{ of 434,112 edges})
\end{equation}
This structural count agrees within 5\% of the global SVD bound (115,223). Since
$\boldsymbol{\delta} \to \lambda\boldsymbol{\delta}$ for any $\lambda \in \mathbb{R}$,
the equivalence class $[\mathbf{W}^*]$ is a \emph{continuous}
$\sim 121{,}100$-dimensional affine subspace of $\mathbb{R}^{434{,}112}$ ---
in contrast to E\&C's discrete equivalence classes of measure zero.  The
biological and experimental consequences of this partition are developed in
\S\ref{subsec:partition}.

%------------------------------------------------------------------------------
\subsection{Identifiable vs.\ duplicate-type edge partition}
\label{subsec:partition}
%------------------------------------------------------------------------------

The equivalence class analysis partitions all 434,112 edges into two
qualitatively different regimes, determined by the local connectivity structure
of the postsynaptic neuron:

\paragraph{Identifiable edges.}
An edge $W_{ij}$ is an \textbf{identifiable edge} if postsynaptic neuron $i$
receives exactly one input of presynaptic type $\alpha = \mathrm{type}(j)$
(i.e.\ $k_{i\alpha} = 1$).  The corresponding column in $\mathbf{H}_i$ is
unique and cannot be redistributed among other columns: $W_{ij}$ is uniquely
fixed by the linear system $\mathbf{H}_i\mathbf{w}_i = \mathbf{b}_i$, so its
equivalence class is a singleton $[\mathbf{w}_{i,j}^*] = \{W_{ij}^*\}$.
Examples of cell types with small in-degree and no duplicate-type inputs include
R1, L3, and Mi1 as postsynaptic targets.

\paragraph{Duplicate-type edges.}
An edge $W_{ij}$ is a \textbf{duplicate-type edge} if neuron $i$ receives
$k_{i\alpha} > 1$ inputs from the same presynaptic type $\alpha$.  The
$k_{i\alpha}$ activity columns in $\mathbf{H}_i$ are strongly correlated ---
each carries the same photoreceptor-driven signal shifted across visual columns
--- contributing effective rank $\approx 1$ instead of $k_{i\alpha}$.  Only
the group sum $\sum_{j':\,\mathrm{type}(j')=\alpha} W_{ij'}$ is constrained by
dynamics; the $k_{i\alpha}-1$ redistributive degrees of freedom lie in
$\ker(\mathbf{H}_i)$.  Group sizes range from $k=2$ for near-isolated connections up to $k \approx 23$
for the most connected types (e.g.\ T4c, $\bar{k}=23$), reflecting the local
column reach of each cell type in the 217-column flyvis connectome.

\paragraph{Biological prediction.}
This partition makes a testable cross-animal prediction about synapse count
conservation.  For \emph{identifiable edges}, $W_{ij}$ is uniquely determined
by the neural dynamics of individual fly $f$: the individual synapse count
$n_{ij}^{(f)}$ should be conserved across flies ($\mathrm{CV}(n_{ij})$ small
across animals).  For \emph{duplicate-type edges}, only the total type-input
$\sum_{j' \in \alpha} W_{ij'}$ is fixed by dynamics: the corresponding
prediction is that \emph{total synaptic input per type per target} is conserved
across animals, not individual synapse counts.  Variability in individual
$n_{ij}^{(f)}$ within the group is unconstrained --- only the group aggregate
is biologically meaningful.

\paragraph{Graded conservation strength.}
The conservation strength is graded by the within-group activity correlation:
near-identical column signals (correlation $\to 1$) leave the group sum as the
\emph{only} constraint, while independent perturbations ($\sigma \uparrow$)
decorrelate same-type activities and collapse duplicate-type groups toward
identifiable edges.  Optogenetic injection of independent per-column stimuli,
or selective ablation of same-type neurons across columns, would increase
$\mathrm{rank}(\mathbf{H}_i)$ and shrink the duplicate-type null space by the
same mechanism as additive noise.

%==============================================================================
\section{Empirical Verification of the Equivalence Class}
%==============================================================================

If the equivalence class analysis is correct, any $\hat{\mathbf{W}} \in
[\mathbf{W}^*]$ should preserve rollout $R^2 \approx 1$ while connectivity $R^2$
can be arbitrarily reduced. We verify this directly.

\subsection{Non-degenerate types: identifiable neurons}

Of 65 cell types, 13 lack duplicate-type groups: R1, L3, Lawf2, Mi1, Mi4, Mi9,
Mi10, Mi11, Mi12, Mi13, T1, Tm5b, Tm9. Their equivalence classes are singletons
($\dim[\mathbf{w}_i^*] = 0$) --- all incoming weights are identifiable. These 13
types are excluded from variant generation.

\subsection{Variant generation from the equivalence class}

For each of the 52 cell types with duplicate-type groups, we generate sum-preserving
perturbations $\boldsymbol{\delta}$ (with $\sum_j\delta_j = 0$, i.e.\
$\boldsymbol{\delta} \in [\mathbf{W}^*]$) at 15 logarithmically-spaced scales
($\lambda = 0.05$ to $8.0$), yielding 960 single-type variants. Mixed-type variants
simultaneously perturb the top 10 cell types by total equivalence class dimension,
generating 1{,}000 additional variants.

\subsection{Results}
\label{subsec:results}

Single-type variants: mean connectivity $R^2 = 0.980 \pm 0.076$, rollout $R^2 =
0.997 \pm 0.045$. Mixed-type variants: $0.839 \pm 0.164$ and $0.997 \pm 0.004$
respectively. The decoupling confirms the structure of $[\mathbf{W}^*]$: large
displacement within the equivalence class leaves dynamics invariant
(Fig.~\ref{fig:null_space_verification}).

At $\lambda = 8.0$ (Table~\ref{tab:lambda_8}), Tm5c achieves connectivity $R^2 =
0.282$ (far from $\mathbf{W}^*$) while rollout $R^2 = 0.997$ --- it lies in a
large-dimensional equivalence class. Tm5Y achieves $R^2 = 1.0$ at the same
perturbation scale, indicating minimal duplicate-type structure in its inputs.

\begin{figure}[h]
\centering
\includegraphics[width=1.0\textwidth]{scatterplot_r2_variants.png}
\caption{Connectivity $R^2$ versus rollout $R^2$ for 1,960 members of the
equivalence class $[\mathbf{W}^*]$. Single-type: 957 variants across 52 cell types
at 15 perturbation scales. Mixed-type: 1,000 variants perturbing top 10 types
simultaneously. The horizontal band at rollout $R^2 \approx 1$ visualizes the
equivalence class directly.}
\label{fig:null_space_verification}
\end{figure}

\begin{table}[h!]
\centering
\footnotesize
\begin{tabular}{lrr@{\quad}lrr@{\quad}lrr}
\toprule
Cell Type & Conn $R^2$ & Rollout $R^2$ & &  &  &  &  &  \\
\midrule
R2 & \cellcolor{corange} 0.9421 & 1.0000 & CT1(M10) & \cellcolor{corange} 0.8774 & 0.9989 & Tm4 & \cellcolor{cgreen} 0.9945 & 1.0000 \\
R3 & \cellcolor{cgreen} 0.9672 & 1.0000 & Mi2 & \cellcolor{corange} 0.9146 & 0.9809 & Tm5Y & \cellcolor{cgreen} 1.0000 & 1.0000 \\
R4 & \cellcolor{corange} 0.7408 & 0.9974 & Mi3 & \cellcolor{corange} 0.9375 & 0.9984 & Tm5a & \cellcolor{cgreen} 0.9694 & 0.9996 \\
R5 & \cellcolor{cgreen} 0.9921 & 0.9999 & Mi14 & \cellcolor{cgreen} 0.9649 & 0.9999 & Tm5c & \cellcolor{cred} 0.2823 & 0.9965 \\
R6 & \cellcolor{cgreen} 0.9995 & 1.0000 & Mi15 & \cellcolor{corange} 0.8710 & 0.9994 & Tm16 & \cellcolor{cgreen} 0.9899 & 1.0000 \\
R7 & \cellcolor{cgreen} 0.9987 & 1.0000 & T2 & \cellcolor{cgreen} 0.9692 & 0.9922 & Tm20 & \cellcolor{cgreen} 0.9966 & 1.0000 \\
R8 & \cellcolor{corange} 0.8993 & 0.9998 & T2a & \cellcolor{cgreen} 0.9688 & 0.9998 & Tm28 & \cellcolor{corange} 0.7526 & 0.9869 \\
L1 & \cellcolor{corange} 0.7275 & 0.9979 & T3 & \cellcolor{cgreen} 0.9664 & 0.9996 & Tm30 & \cellcolor{corange} 0.5658 & 0.9982 \\
L2 & \cellcolor{corange} 0.6171 & 0.9980 & T4a & \cellcolor{cgreen} 0.9668 & 0.9983 & TmY3 & \cellcolor{cgreen} 0.9983 & 1.0000 \\
L4 & \cellcolor{cgreen} 0.9961 & 0.9983 & T4b & \cellcolor{cgreen} 0.9904 & 0.9985 & TmY4 & \cellcolor{cgreen} 0.9931 & 1.0000 \\
L5 & \cellcolor{corange} 0.7332 & 0.9944 & T4c & \cellcolor{cgreen} 0.9875 & 0.9990 & TmY5a & \cellcolor{corange} 0.9409 & 0.9999 \\
Lawf1 & \cellcolor{cgreen} 0.9967 & 1.0000 & T4d & \cellcolor{cgreen} 0.9778 & 0.9988 & TmY9 & \cellcolor{cgreen} 0.9517 & 0.9992 \\
Am & \cellcolor{cgreen} 0.9512 & 0.9999 & T5a & \cellcolor{cgreen} 0.9975 & 0.9999 & TmY10 & \cellcolor{cgreen} 0.9723 & 0.9616 \\
C2 & \cellcolor{corange} 0.9270 & 0.9996 & T5b & \cellcolor{cgreen} 0.9977 & 1.0000 & TmY13 & \cellcolor{cgreen} 0.9942 & 1.0000 \\
\bottomrule
\end{tabular}
\caption{Single-type variants at extreme scale $\lambda=8.0$ for 51 cell types with
duplicate-type groups. Connectivity $R^2$ measures distance from $\mathbf{W}^*$
within $[\mathbf{W}^*]$; rollout $R^2$ confirms membership in the equivalence
class. Green ($R^2>0.95$), orange ($0.50<R^2\leq 0.95$), red ($R^2\leq 0.50$).}
\label{tab:lambda_8}
\end{table}

%------------------------------------------------------------------------------
\subsection{Cross-seed GNN ensemble: sampling the equivalence class}
\label{subsec:cv-ensemble}
%------------------------------------------------------------------------------

Each gradient-descent run from a different random seed converges to a
\emph{different} member of $[\mathbf{W}^*]$.  Running $K=5$ independent
training seeds on the same data (same simulation seed, same architecture)
produces five connectivity matrices $\hat{\mathbf{W}}^{(1)},\ldots,
\hat{\mathbf{W}}^{(5)}$ that all achieve similar rollout performance but
vary substantially in their edge values.  Stacking them into a
$5\times 434{,}112$ matrix and computing its SVD reveals the structure of
the manifold region explored by gradient descent.

\paragraph{Per-seed performance.}

\begin{center}
\begin{tabular}{lcc}
\toprule
Seed & Connectivity $R^2$ & Rollout Pearson $r$ \\
\midrule
cv00 & 0.934 & $>0.99$ \\
cv01 & 0.969 & $>0.99$ \\
cv02 & 0.968 & $>0.99$ \\
cv03 & 0.967 & $>0.99$ \\
cv04 & 0.972 & $>0.99$ \\
\midrule
\textbf{Consensus mean} & \textbf{0.976} & $>0.99$ \\
Consensus median & 0.978 & $>0.99$ \\
\bottomrule
\end{tabular}
\end{center}

The consensus estimate $\bar{\mathbf{W}} = K^{-1}\sum_k \hat{\mathbf{W}}^{(k)}$
improves over the best individual seed by $+0.004$ in $R^2$.  Averaging
cancels the null-space components---which are random across seeds---while
preserving the identifiable component, consistent with $\bar{\mathbf{W}}$
being the minimum-variance unbiased estimator within the identifiable subspace.

\paragraph{Empirical null-space directions.}

Centering the seed stack gives $\tilde{\mathbf{W}} = \hat{\mathbf{W}} -
\bar{\mathbf{W}}$ (shape $5\times 434{,}112$).  Its SVD yields the principal
directions of inter-seed variation:
\begin{center}
\begin{tabular}{ccc}
\toprule
Component & Singular value & \% variance explained \\
\midrule
1 & 23.2 & 58.1\% \\
2 & 13.3 & 19.2\% \\
3 & 11.4 & 13.9\% \\
4 &  9.0 &  8.7\% \\
5 & $\approx 0$ & $\approx 0\%$ \\
\bottomrule
\end{tabular}
\end{center}
Five seeds span a $\leq 4$-dimensional subspace of the
$\sim 121{,}100$-dimensional null space.  The right singular vectors
$\mathbf{v}_k \in \mathbb{R}^{434{,}112}$ are empirical basis vectors of
$\ker(\mathbf{H})$; visualizing their edge-weight support identifies the cell
types and projection groups most responsible for residual ambiguity.

\paragraph{Measuring identifiability at the edge level.}

The identifiable/duplicate-type classification is fundamentally a property of
each \emph{edge}, not each cell type: an edge $W_{ij}$ depends on the
combination of postsynaptic neuron $i$ (which sets how many same-type inputs
exist, $k_{i\alpha}$) and presynaptic type $\alpha = \mathrm{type}(j)$.
Formally, the degeneracy of a given edge is
\begin{equation}
d_{ij} = k_{i,\alpha} - 1 \;\geq 0
\end{equation}
($d_{ij}=0$: identifiable; $d_{ij}>0$: $d_{ij}$ free redistributive
directions).  The natural structural metric is therefore computed per
$(i, \alpha)$ group, or aggregated by presynaptic type $\alpha$:
\begin{equation}
\bar{k}_\alpha = \frac{1}{|E_\alpha|}
\sum_{(i,j)\in E_\alpha} k_{i,\alpha}
\end{equation}
where $E_\alpha$ is the set of all edges with $\mathrm{type}(j) = \alpha$.
$\bar{k}_\alpha = 1$ means all outgoing edges of type $\alpha$ land on targets
with exactly one input of that type (fully identifiable); $\bar{k}_\alpha \gg 1$
means most outgoing edges belong to large degenerate groups.

As the empirical counterpart we compute the per-type \textbf{seed SNR}:
\begin{equation}
\mathrm{SNR}(\alpha) = \left\langle\frac{|\bar{W}_e|}{\sigma_e}\right\rangle_{e\in E_\alpha}
\end{equation}
where $\bar{W}_e = K^{-1}\sum_k \hat{W}_e^{(k)}$ and $\sigma_e =
\mathrm{std}_k(\hat{W}_e^{(k)})$ are the cross-seed mean and standard
deviation.  High $\mathrm{SNR}(\alpha)$ confirms that seeds agree on
those weights; low $\mathrm{SNR}(\alpha)$ means null-space redistribution
dominates --- the seeds found the same group sums but placed the individual
weights differently.

\begin{table}[h!]
\centering
\footnotesize
\begin{tabular}{lrc|lrc}
\toprule
Cell type & $\bar{k}_\alpha$ & SNR &
Cell type & $\bar{k}_\alpha$ & SNR \\
\midrule
\multicolumn{3}{c|}{\textit{Seeds agree (high SNR)}} &
\multicolumn{3}{c}{\textit{Seeds disagree (low SNR)}} \\
\midrule
R8       & 1.6 & 95  & L2     & 4.5 & 0.5 \\
CT1(Lo1) & 7.8 & 93  & L1     & 9.2 & 0.5 \\
Mi2      & 3.7 & 83  & T1     & 1.0 & 0.5 \\
Tm5c     & 1.5 & 80  & Am     & 1.0 & 0.6 \\
Mi4      & 6.1 & 61  & Lawf1  & 1.0 & 0.6 \\
T3       & 2.3 & 54  & Tm28   & 4.6 & 0.7 \\
Tm9      & 3.2 & 54  & TmY3   & 8.1 & 3.5 \\
Tm20     & 4.8 & 45  & Lawf2  & 3.8 & 3.7 \\
R7       & 2.4 & 42  &        &     &     \\
\bottomrule
\end{tabular}
\caption{Per-type structural multiplicity $\bar{k}_\alpha$ (mean number of same-type
inputs at the target neuron, computed from the flyvis connectome) and empirical seed SNR
$= \langle|\bar{W}_e|/\sigma_e\rangle$ aggregated over outgoing edges of each source
type $\alpha$ across five CV training seeds.  Both quantities measure identifiability
of \emph{outgoing} connections from type $\alpha$; compare with
Table~\ref{tab:lambda_8}, which measures degeneracy of \emph{incoming} weights.
Structural and empirical measures are not always aligned: L1 and L2 show
high $\bar{k}_\alpha$ and low SNR (null-space degeneracy from many same-type projections),
while T1, Am, and Lawf1 show $\bar{k}_\alpha=1$ (structurally identifiable) but also
low SNR, consistent with near-zero absolute weights where any seed variation
dominates the ratio.}
\label{tab:snr}
\end{table}

The SNR analysis reveals two distinct mechanisms for low seed agreement.
\emph{Null-space degeneracy}: L1 ($\bar{k}=9.2$) and L2 ($\bar{k}=4.5$) project
heavily into large duplicate-type groups at their targets; seeds redistribute
weight within the group sum, yielding $\mathrm{SNR}\leq 0.5$.
\emph{Near-zero signal}: T1, Am, and Lawf1 have $\bar{k}=1$ (all outgoing edges
structurally identifiable) yet also $\mathrm{SNR}<0.7$, indicating that their
absolute synaptic weights are near zero in the learned model, so any inter-seed
fluctuation swamps the signal.  By contrast, R8, CT1(Lo1), and Mi4 recover
consistently (SNR $\geq 60$) despite varying $\bar{k}$, showing that strong
synaptic weights dominate even moderate null-space noise.  Complete per-type
values are provided in the supplementary CSV ({\tt scripts/type\_identifiability.csv}).

\paragraph{Edge-level validation of the partition.}
The identifiable/duplicate-type classification is properly a property of each
\emph{edge} $W_{ij}$, determined by the group size $k_{i,\alpha}$ of the
specific (postsynaptic neuron $i$, presynaptic type $\alpha$) pair --- not of
cell types as a whole.  Of the 434,112 edges, 61,673 (14.2\%) are identifiable
($k_{i,\alpha}=1$) and 372,439 (85.8\%) are duplicate-type ($k_{i,\alpha}>1$).
Most source types have \emph{both} kinds of outgoing edges: 52 of 65 types
project to targets with mixed $k$ values.  Eleven types send only identifiable
outputs ($f_\text{ident}=1$: Am, R1--R6, Mi3, Lawf1, T1, Tm5Y), while seven
types are near-pure duplicate ($f_\text{ident}<0.05$, e.g.\ Mi2, Mi13, T3,
TmY15).

A direct edge-level test of the theory confirms the correct direction but with
a small effect: identifiable edges have higher cross-seed SNR than duplicate-type
edges with AUROC $= P(\mathrm{SNR}_{k=1} > \mathrm{SNR}_{k>1}) = 0.511$
(Mann-Whitney $p < 10^{-17}$, Cohen's $d = 0.08$ on log-SNR).  The marginal
separation arises because $\mathrm{SNR} = |\bar{W}_e|/\sigma_e$ conflates two
factors: the \emph{null-space noise} $\sigma_e$ (elevated for $k>1$ edges by
the $k-1$ redistributive degrees of freedom) and the \emph{weight magnitude}
$|\bar{W}_e|$ (an independent biological property).  A within-type controlled
comparison eliminates this confound: for the 52 mixed source types, identifiable
outgoing edges have higher mean SNR than duplicate-type outgoing edges in 41 of
52 cases (79\%).  The 11 exceptions arise either when identifiable edges of a
type carry near-zero weights (e.g.\ T5d: $\mathrm{SNR}_{k=1}=0.6$ vs
$\mathrm{SNR}_{k>1}=24$), or when a duplicate-type group has exceptionally
strong mean weights that dominate the null-space variance (e.g.\ R8:
$\mathrm{SNR}_{k=1}=63$ vs $\mathrm{SNR}_{k>1}=143$).
Edge-level results are tabulated in {\tt scripts/edge\_partition\_snr.csv}.

%------------------------------------------------------------------------------
\subsection{Minimum-support member of the equivalence class}
\label{subsec:sparse}
%------------------------------------------------------------------------------

Within $[\mathbf{W}^*]$, the \emph{minimum-support} (maximally sparse)
representative is obtained by collapsing each degenerate group onto a single
edge.  For each $(i,\alpha)$ group of $k_{i\alpha}$ edges sharing the same
postsynaptic neuron $i$ and presynaptic type $\alpha$, we select a
representative edge $e^* = \arg\max_{e \in \text{group}} |W_e^*|$ and set:
\begin{equation}
\hat{W}_e =
\begin{cases}
W_{\text{rep}} & e = e^* \\
0 & \text{otherwise}
\end{cases}
\end{equation}
\emph{Sum-preserving variant}: $W_{\text{rep}} = \sum_{e \in \text{group}} W_e^*$ (exact null-space member for identical group activities). \emph{Calibrated variant}: $W_{\text{rep}}$ is the scalar minimising the per-group dynamics residual over observed voltage trajectories,
\begin{equation}
W_{\text{rep}}^* = \frac{\sum_t h_{e^*}(t)\,\sum_{e\in\text{group}} W_e^*\,h_e(t)}
                        {\sum_t h_{e^*}(t)^2}
\end{equation}
which is a single least-squares problem per group (one degree of freedom).

Applying this to all 64,279 degenerate groups zeros
$\sum_{i,\alpha}(k_{i\alpha}-1)=308{,}160$ edges --- 71\% of the 434,112
total.  Results:

\begin{center}
\begin{tabular}{lrrrr}
\toprule
Method & Edges zeroed & Conn $R^2$ & Rollout $R^2$ & Rollout Pearson $r$ \\
\midrule
Sum-preserving & 308,160 (72.5\%) & 0.105 & 0.944 & 0.974 \\
Calibrated     & 308,160 (72.5\%) & 0.394 & 0.983 & \textbf{0.991} \\
\bottomrule
\end{tabular}
\end{center}

The calibrated method achieves rollout Pearson $r > 0.99$ after zeroing 72.5\%
of edges, confirming approximate membership in $[\mathbf{W}^*]$ far from the
ground truth $\mathbf{W}^*$.  The residual rollout gap (sum-preserving:
$R^2=0.944$; calibrated: $R^2=0.983$) quantifies the \emph{approximation
error} of the structural null space: same-type group members have correlated
but not identical activities, so exact degeneracy is replaced by approximate
degeneracy and a non-trivial representative weight is needed.

The low connectivity $R^2$ of the calibrated sparse solution (0.394) is
expected by construction: each representative absorbs its entire group's
weight, making it appear $k$-fold larger than the corresponding GT edge.
The dynamics are preserved; the individual edge values are not.

%------------------------------------------------------------------------------
\subsection{Pseudoinverse recovery}
\label{subsec:pseudoinverse}
%------------------------------------------------------------------------------

We compare three linear recovery methods. \emph{Truncated SVD}: pseudoinverse
$\mathbf{w}_i^+ = \mathbf{H}_i^+\mathbf{b}_i$ (threshold $10^{-12}$). \emph{Ridge}:
Tikhonov regularization $\lambda=1$. \emph{CV-ridge}: per-neuron $\lambda$ by
time-split cross-validation.

\begin{center}
\begin{tabular}{llrrr}
\toprule
Method & Condition & Conn $R^2$ & RMSE$(T)$ & Rollout $R^2(T)$ \\
\midrule
truncated SVD & $\sigma = 0$    & 0.988 & $9.7\times10^{-5}$ & 1.000 \\
              & $\sigma = 0.05$ & 0.991 & $1.0\times10^{-6}$ & 1.000 \\
              & $\sigma = 0.5$  & 1.000 & $9.2\times10^{-7}$ & 1.000 \\
\midrule
ridge ($\lambda=1$) & $\sigma=0$    & 0.879 & $7.1\times10^{-3}$ & 1.000 \\
                    & $\sigma=0.05$ & 0.989 & $1.9\times10^{-3}$ & 1.000 \\
                    & $\sigma=0.5$  & 0.999 & $6.4\times10^{-4}$ & 1.000 \\
\midrule
CV-ridge & $\sigma=0$    & 0.991 & $2.2\times10^{-6}$ & 1.000 \\
         & $\sigma=0.05$ & 0.991 & $8.5\times10^{-7}$ & 1.000 \\
         & $\sigma=0.5$  & 1.000 & $9.4\times10^{-7}$ & 1.000 \\
\bottomrule
\end{tabular}
\end{center}

Truncated SVD and CV-ridge achieve near-identical results ($R^2 \geq 0.988$):
the system is well-conditioned within the identifiable subspace. Fixed ridge
over-regularizes at $\sigma=0$ (0.879), shrinking identifiable directions
unnecessarily.

Near-perfect recovery at $\sigma=0$ (0.988) clarifies the structure of
$[\mathbf{W}^*]$: the non-identifiable directions have small but nonzero singular
values. The equivalence class is \emph{approximate}, not exact --- duplicate-type
neuron activities are highly correlated but not identical, and the pseudoinverse
exploits residual variation to select $\mathbf{w}_i^*$ from $[\mathbf{w}_i^*]$.

%==============================================================================
\section{Discussion}
%==============================================================================

\subsection{Continuous vs.\ discrete equivalence classes}

In teacher-student MLP recovery (Martinelli et al., 2024), the equivalence class
of zero-loss students is a \emph{discrete} finite set arising from permutation,
sign-flip, and scaling symmetries. Here the equivalence class $[\mathbf{W}^*]$ is
a \emph{continuous} $\sim 121{,}100$-dimensional affine manifold, arising from the
biological duplicate-type symmetry of the columnar connectome. This makes the
inverse problem structurally harder: the manifold has positive measure in
$\mathbb{R}^{434{,}112}$, so any unregularized estimator converges to an arbitrary
point in $[\mathbf{W}^*]$ rather than to $\mathbf{W}^*$ itself.

\subsection{The GNN as a misspecified teacher-student setup}

Our framework is a misspecified teacher-student setup: the teacher is the physical
ODE simulator (Eq.~\ref{Eq:ODE}) and the student is the GNN:
\begin{equation}
\frac{\widehat{dv}_i(t)}{dt}
= f_\theta\!\Big(v_i(t),\mathbf{a}_i,
\sum_{j\in\mathcal{N}_i}\hat{W}_{ij}\,g_\phi(v_j(t),\mathbf{a}_j)^2,
I_i(t)\Big)
\end{equation}
where $f_\theta,g_\phi$ are MLPs, $\hat{W}_{ij}$ are learnable edge weights, and
$\mathbf{a}_i \in \mathbb{R}^2$ are learnable per-neuron embeddings. The student
must simultaneously learn the dynamics \emph{and} select a point from
$[\mathbf{W}^*]$. Unlike E\&C where the student architecture matches the teacher
exactly, here the student approximates a continuous ODE with discrete MLP
layers --- a fundamental mismatch. Training minimizes:
\begin{equation}
\begin{aligned}
\mathcal{L} &=
\|\widehat{\mathbf{y}}-\mathbf{y}\|_2
+ \lambda_0\|\theta\|_1 + \lambda_1\|\phi\|_1 + \lambda_2\|\hat{\mathbf{W}}\|_1 \\
&\quad + \gamma_0\|\theta\|_2 + \gamma_1\|\phi\|_2 \\
&\quad + \mu_0\|\operatorname{ReLU}(-\partial_v g_\phi)\|_2
       + \mu_1\|g_\phi(v_\star,\mathbf{a})-v_\star\|_2
\end{aligned}
\end{equation}
The $\ell_1$ penalty on $\hat{\mathbf{W}}$ acts as a regularizer that selects a
sparse point from $[\mathbf{W}^*]$, analogously to E\&C's clustering step
selecting the teacher neuron from the set of duplicates. The GNN's message-passing
inductive bias further constrains which member of $[\mathbf{W}^*]$ is found.

The gap between GNN ($R^2=0.899$, $\sigma=0$) and pseudoinverse ($R^2=0.988$)
quantifies the cost of \emph{dynamic misspecification}, not the size of
$[\mathbf{W}^*]$: the pseudoinverse has exact access to $\tau_i$, $V_i^{\text{rest}}$,
and $\dot{v}_i$; the GNN must approximate all three.

\subsection{The noise level $\sigma$ as an ease-of-identifiability axis}

In E\&C, the ratio $r/d_{\text{in}}$ defines an ease-of-identifiability axis:
low $r/d_{\text{in}}$ means orthogonal teacher hyperplanes and easy recovery; high
$r/d_{\text{in}}$ means near-parallel hyperplanes, a rough loss landscape, and
hard recovery. Here the noise level $\sigma$ plays the identical role:

\begin{equation}
\sigma \uparrow \quad\Rightarrow\quad
\text{rank}(\mathbf{H}_i) \uparrow \quad\Rightarrow\quad
\dim[\mathbf{w}_i^*] \downarrow \quad\Rightarrow\quad
R^2_{\mathbf{W}} \uparrow
\end{equation}

\begin{center}
\begin{tabular}{lrrrr}
\toprule
Noise & Activity rank $r$ & $\dim[\mathbf{W}^*]$ & \% identifiable & GNN $R^2_{\mathbf{W}}$ \\
\midrule
$\sigma=0$    & 45  & 115,223 & 73.5\%  & $0.899\pm0.030$ \\
$\sigma=0.05$ & 134 & 10,227  & 97.6\%  & $0.962\pm0.012$ \\
$\sigma=0.5$  & 781 & 0       & 100.0\% & $0.997\pm0.000$ \\
\bottomrule
\end{tabular}
\end{center}

Independent noise realizations break the duplicate-type symmetry: same-type neurons
now receive distinct perturbations, so their activity columns in $\mathbf{H}_i$
become linearly independent, collapsing $[\mathbf{w}_i^*]$ to a singleton. At
$\sigma=0.5$, rank $781 \gg d_i^{\max} = 208$, so every neuron's equivalence class
is a singleton and $R^2_{\mathbf{W}} = 0.997$.

Experimentally, the same collapse could be achieved by optogenetic injection of
independent perturbations per neuron (e.g., distinct light patterns per L1 across
columns), or selective ablation of same-type neurons across columns --- both
increase $\text{rank}(\mathbf{H}_i)$ and shrink $[\mathbf{w}_i^*]$ by the same
mechanism as noise.

%==============================================================================
\section{Summary}
%==============================================================================

The equivalence class $[\mathbf{W}^*]$ of zero-rollout-loss connectomes is a
continuous $\sim 121{,}100$-dimensional affine manifold, arising from
duplicate-type neuron groups in the columnar flyvis connectome ($\sim 28\%$ of
434,112 edges). This is the biological analogue of overparameterisation symmetry
in teacher-student MLP recovery, but continuous rather than discrete. Two
independent estimates of $\dim[\mathbf{W}^*]$ --- global SVD (115,223) and
structural counting ($\sim 121{,}100$) --- agree within 5\%. Empirically, 1,960
variants sampled from $[\mathbf{W}^*]$ confirm the decoupling: connectivity $R^2$
ranges 0.28--1.0 while rollout $R^2 > 0.96$.

\paragraph{Multiple GNN seeds sample $[\mathbf{W}^*]$.}
Five independent training seeds converge to distinct members of $[\mathbf{W}^*]$,
yielding connectivity $R^2 \in [0.934, 0.972]$ at fixed rollout quality.  Their
ensemble mean improves to $R^2=0.976$, as null-space components (random across
seeds) cancel while the identifiable component is preserved.  SVD of the
inter-seed variation reveals a $\leq 4$-dimensional explored subspace of the
$121{,}100$-dimensional null space, with the dominant direction explaining 58\%
of inter-seed variance.  Per-type seed SNR analysis (Table~\ref{tab:snr}) identifies two mechanisms
of low identifiability: null-space degeneracy (L1, L2; $\bar{k} \geq 4.5$,
SNR $\leq 0.5$) and near-zero absolute weights (T1, Am, Lawf1; $\bar{k}=1$
but SNR $< 0.7$).  Types with strong synaptic weights recover consistently
regardless of $\bar{k}$: R8 and CT1(Lo1) achieve SNR $> 90$.

\paragraph{Minimum-support extreme of $[\mathbf{W}^*]$.}
Collapsing each degenerate $(i,\alpha)$ group to a single calibrated
representative zeros 308,160 edges (72.5\%) while maintaining rollout Pearson
$r = 0.991$.  This is the minimum-support member of $[\mathbf{W}^*]$:
the maximum-zeros point on the null-space manifold that still preserves dynamics.
The 308,160 zeros correspond exactly to the structural null dimension of the
64,279 degenerate groups.

Recovery of $\mathbf{W}^*$ from $[\mathbf{W}^*]$ is nonetheless possible because
the equivalence class is \emph{approximate}: duplicate-type activities are
correlated but not identical, and the pseudoinverse selects $\mathbf{W}^*$ with
$R^2 = 0.988$ even at $\sigma=0$. The noise level $\sigma$ defines an
ease-of-identifiability axis analogous to $r/d_{\text{in}}$ in E\&C, collapsing
$[\mathbf{W}^*]$ monotonically: rank $45 \to 781$, GNN $R^2$ $0.899 \to 0.997$.
The residual GNN gap from the pseudoinverse ($0.988$) quantifies dynamic
misspecification, not the intrinsic size of $[\mathbf{W}^*]$.

\end{document}
